\section{A common closed specialization of all quadratic pencils}
\label{sec:closed-pencil-v151}

We give an orbit-closure consequence directly on the pencil side.
Unlike restricting an equivalence to a previously specified locus, the
following argument constructs the specialization.  It works for every
pencil, including pencils with singular generic member.  Its mechanism
is elementary projective representation theory; the additional failure
algebra statement uses the all-pencil, exact-multiplicity construction.

\begin{theorem}[Explicit common closed orbit]
\label{thm:closed-pencil-v151}
Let \(\dim V=n\geq3\).  The action of \(\PGL(V)\) on
\(\mathscr P=\Gr(2,\Sym^2V)\) has a unique closed orbit, namely
\begin{equation}\label{eq:flag-orbit-v151}
 \mathscr C=\{\ell H:\ell\subset H\subset V,\
                      \dim\ell=1,\ \dim H=2\}.
\end{equation}
It is the flag variety \(\operatorname{Fl}(1,2;V)\), of dimension
\(2n-3\).  More precisely, after a fixed change of coordinates every
pencil has a degeneration along the one-parameter subgroup
\[
                  \lambda(t)=\diag(1,t,t^2,\ldots,t^2)
\]
to \(\langle x_1^2,x_1x_2\rangle\).  The pencil family is a
subbundle on the whole affine line.
\end{theorem}
\begin{proof}
Take independent forms \(q_0,q_1\).  The rational function
\(q_1/q_0\) on \(\Pj(V^*)\) is nonconstant.  In characteristic
zero its differential is not identically zero: on an affine coordinate
chart the common kernel of the coordinate derivations of the rational
function field is \(\C\).  Choose a point \([v]\) where \(q_0(v)\ne0\)
and this differential is nonzero.  Replace \(q_1\) by
\(q_1-(q_1(v)/q_0(v))q_0\).  Then \(q_1(v)=0\), whereas its
differential has a nonzero value on some vector \(w\) independent of
\(v\).  Euler's identity shows that its value on the radial vector
\(v\) is zero.  Extending \(v,w\) to a basis of \(V^*\) gives dual
coordinates in which
\[
 [x_1^2]q_0=a\ne0,\qquad [x_1^2]q_1=0,
                                  \qquad [x_1x_2]q_1=b\ne0.
\]
Here brackets denote a coefficient.

Apply \(\lambda(t)\) to the variables.  The forms
\[
 Q_0(t)=q_0(x_1,tx_2,t^2x_3,\ldots,t^2x_n),\qquad
 Q_1(t)=t^{-1}q_1(x_1,tx_2,t^2x_3,\ldots,t^2x_n)
\]
have polynomial coefficients: the only possible weight-zero term in
\(q_1\) was removed.  Their limits are \(a x_1^2\) and
\(b x_1x_2\).  For \(t\ne0\) they span a coordinate transform of
\(R\); at zero they are independent.  Thus their coefficient matrix
has rank two everywhere, defines a subbundle, and supplies the asserted
morphism from \(\A^1\) to \(\mathscr P\).

The orbit of \(\langle x_1^2,x_1x_2\rangle\) is exactly
\(\mathscr C\).  A pencil of the form \(\ell H\) recovers \(\ell\)
as its gcd and \(H\) by division.  These are scheme-theoretic inverse
operations on the exact-gcd-one-degree stratum by
Corollary~\ref{cor:canonical-exact-v151}, applied to \(g=h=1,r=2\).
Under that inverse the condition \(\ell\subset H\) is the usual
closed flag incidence.  The flag-to-pencil map is therefore a locally
closed immersion; since the flag variety is projective it is a closed
immersion.  It identifies \(\mathscr C\) with the flag variety, and
the dimension is \((n-1)+(n-2)\).  Every pencil orbit closure contains
\(\mathscr C\), by the constructed family and equivariance.  Any
closed orbit must consequently be \(\mathscr C\), proving uniqueness.
\end{proof}

\begin{corollary}[Flat closed specialization of failure algebras]
\label{cor:closed-failure-v151}
For every quadratic pencil \(R\) there is a finite locally free
family over \(\A^1\), of rank
\(\binom{n^2+d}{d}-\binom N2\), whose nonzero fibres are isomorphic
to \(\mathfrak A_R^{[d]}\) and whose zero fibre is
\(\mathfrak A_{\langle x_1^2,x_1x_2\rangle}^{[d]}\).
The intrinsic first relations throughout the family have exact gcd
\(\delta^{n-1}\).  In a fixed determinant frame, the two ordered
copies of the flag orbit form the unique closed \(G\)-orbit in the
pencil relation locus.  Equivalently the effective pencil failure
stack has the common closed specialization described in
Theorem~\ref{thm:closed-pencil-v151}.
\end{corollary}
\begin{proof}
Let \(R_t\) be the polynomial subbundle just constructed.  Its
Pluecker line gives, under the fixed exterior-duality identification,
a line bundle \(L_t^0\subset\mathbb S_\mu V^*\otimes\OO_{\A^1}\).
Set
\[
 \mathcal K_t=\delta^{n-1}(L_t^0\otimes\mathbb S_\mu U),\qquad
 \mathcal A_t=\Sym(E\otimes\OO_{\A^1})/
                       ((\mathcal K_t)+\mathfrak m^{d+1}).
\]
The first is a rank-\(M\) subbundle.  The second is finite locally
free by its graded direct-sum presentation, with the claimed rank.
Lemma~\ref{lem:primitive-pencil} gives the exact gcd at every fibre,
including zero; in fact this family lies scheme-theoretically in the
fixed-divisor exact stratum.  Equivariance and the local inverse
identify its nonzero fibres as asserted.  In a determinant frame the
pencil locus is \(G\times^{G^\circ}\mathscr P\).  The previous
closed-orbit calculation therefore gives the assertion for \(G\).
This is a statement inside that locus; it does not assert closedness
under arbitrary degenerations of the determinant form in the full
relation Grassmannian.
\end{proof}

The closed specialization is very coarse compared with the unmarked
local inverse: many pairwise nonisomorphic finite algebras have this
same limiting fibre.  The full first relation, not its common limit,
retains the original pencil.  The explicit family also differs from
transport of spectral rank loci: it crosses from arbitrary pencils to
a singular common-factor pencil without a regularity hypothesis.
Theorem~\ref{thm:boundary-normalization-v151}, on the other hand,
describes divisor collisions in the larger natural first-relation
boundary, outside the exact pencil image.  These are two distinct
geometric statements, and neither is a classification of every
congruence-orbit adjacency of quadratic pencils.
